
\documentclass{article}
\usepackage{colortbl}
\usepackage{makecell}
\usepackage{multirow}
\usepackage{supertabular}

\begin{document}

\newcounter{utterance}

\centering \large Interaction Transcript for game `cladder', experiment `full\_v1.5\_default', episode 2271 with score{-}Qwen3.5{-}9B{-}sp.
\vspace{24pt}

{ \footnotesize  \setcounter{utterance}{1}
\setlength{\tabcolsep}{0pt}
\begin{supertabular}{c@{$\;$}|p{.15\linewidth}@{}p{.15\linewidth}p{.15\linewidth}p{.15\linewidth}p{.15\linewidth}p{.15\linewidth}}
    \# & \multicolumn{2}{c}{Player} && \multicolumn{2}{c}{Game Master} \\
    \hline

    \theutterance \stepcounter{utterance}  
    & & & \multicolumn{4}{p{0.6\linewidth}}{
        \cellcolor[rgb]{0.9,0.9,0.9}{
            \makecell[{{p{\linewidth}}}]{
                \texttt{\tiny{[P1$\langle$GM]}}
                \texttt{You are an expert in causal inference. The following question is not a typical commonsense query, but rather a meticulously designed question created by a professor specializing in causal inference, intended to assess the students' mastery of the course content.} \\
\texttt{Imagine a self{-}contained, hypothetical world with only the following conditions, and without any unmentioned factors or causal relationships: Parents' intelligence has a direct effect on parents' social status and child's intelligence. Other unobserved factors has a direct effect on parents' social status and child's intelligence. Parents' social status has a direct effect on child's intelligence. Other unobserved factors is unobserved. The overall probability of intelligent parents is 15\%. The probability of unintelligent parents and intelligent child is 54\%. The probability of intelligent parents and intelligent child is 4\%. Is the chance of intelligent child smaller when observing intelligent parents?} \\
\texttt{Start your answer with 'Yes' or 'No', followed by additional reasoning or evidence to support your explanation.} \\
            }
        }
    }
     \\ \\

    \theutterance \stepcounter{utterance}  
    & \multicolumn{4}{p{0.6\linewidth}}{
        \cellcolor[rgb]{0.561,0.737,0.561}{
            \makecell[{{p{\linewidth}}}]{
                \texttt{\tiny{[P1$\rangle$GM]}}
                \texttt{No} \\
\\ 
\texttt{**Reasoning:**} \\
\\ 
\texttt{To determine whether the chance of an intelligent child is smaller when observing intelligent parents, we must compare the conditional probability \$P(\textbackslash{}text\{Intelligent Child\} \textbackslash{}mid \textbackslash{}text\{Intelligent Parents\})\$ with the marginal probability \$P(\textbackslash{}text\{Intelligent Child\})\$.} \\
\\ 
\texttt{1.  **Calculate the Marginal Probability of an Intelligent Child (\$P(C)\$):**} \\
\texttt{    We are given the following probabilities:} \\
\texttt{    *   \$P(P\_\{int\}) = 0.15\$ (Probability of intelligent parents)} \\
\texttt{    *   \$P(P\_\{int\} \textbackslash{}cap C\_\{int\}) = 0.04\$ (Probability of intelligent parents AND intelligent child)} \\
\texttt{    *   \$P(P\_\{unint\} \textbackslash{}cap C\_\{int\}) = 0.54\$ (Probability of unintelligent parents AND intelligent child)} \\
\\ 
\texttt{    The total probability of an intelligent child is the sum of the probabilities of an intelligent child occurring with intelligent parents and with unintelligent parents:} \\
\texttt{    \$\$P(C\_\{int\}) = P(P\_\{int\} \textbackslash{}cap C\_\{int\}) + P(P\_\{unint\} \textbackslash{}cap C\_\{int\})\$\$} \\
\texttt{    \$\$P(C\_\{int\}) = 0.04 + 0.54 = 0.58\$\$} \\
\\ 
\texttt{2.  **Calculate the Conditional Probability of an Intelligent Child given Intelligent Parents (\$P(C\_\{int\} \textbackslash{}mid P\_\{int\})\$):**} \\
            }
        }
    }
    & & \\ \\

    \theutterance \stepcounter{utterance}  
    & & & \multicolumn{2}{p{0.3\linewidth}}{
        \cellcolor[rgb]{0.9,0.9,0.9}{
            \makecell[{{p{\linewidth}}}]{
                \texttt{\tiny{[GM$|$GM]}}
                \texttt{No} \\
            }
        }
    }
    & & \\ \\

    \theutterance \stepcounter{utterance}  
    & & & \multicolumn{2}{p{0.3\linewidth}}{
        \cellcolor[rgb]{0.9,0.9,0.9}{
            \makecell[{{p{\linewidth}}}]{
                \texttt{\tiny{[GM$|$GM]}}
                \texttt{yes} \\
            }
        }
    }
    & & \\ \\

    \theutterance \stepcounter{utterance}  
    & & & \multicolumn{2}{p{0.3\linewidth}}{
        \cellcolor[rgb]{0.9,0.9,0.9}{
            \makecell[{{p{\linewidth}}}]{
                \texttt{\tiny{[GM$|$GM]}}
                \texttt{game\_result = LOSE} \\
            }
        }
    }
    & & \\ \\

\end{supertabular}
}

\end{document}
